\chapter{Górne twierdzenie Löwenheima--Skolema i paradoks Skolema}
\label{ch:upward-ls}

\section{Dlaczego można powiększać model bez zmiany jego zdań?}
Dolne twierdzenie Löwenheima--Skolema
(Twierdzenie~\ref{thm:downward-ls}) zapewnia małe
podstruktury elementarne. Teraz skonstruujemy rozszerzenia
elementarne nieskończonej struktury o dowolnie dużej mocy.
Nie wystarczy dodać elementów: prawdziwość wszystkich
formuł z parametrami ze starej struktury musi pozostać
niezmieniona. W tym celu nazwiemy wszystkie jej elementy
i zachowamy pełny zbiór zdań, które są o nich prawdziwe.

\section{Diagram elementarny modelu}
\begin{definicja}[Rozszerzenie języka o nazwy elementów]
\label{def:ls-expanded-language}
Niech $\mathcal L$ będzie językiem pierwszego rzędu,
a $\mathcal M$ strukturą dla $\mathcal L$
o niepustym uniwersum $M$.
Do $\mathcal L$ dodajemy nowe, różne symbole
stałych $c_a$ dla wszystkich $a\in M$;
otrzymany język oznaczamy $\mathcal L(M)$.
Przez $\mathcal M_M$ oznaczamy rozszerzenie
struktury $\mathcal M$ do języka $\mathcal L(M)$
spełniające $c_a^{\mathcal M_M}=a$.
\end{definicja}

\begin{definicja}[Diagram elementarny]
\label{def:ls-elementary-diagram}
Dla struktury $\mathcal M$ z poprzedniej definicji
jej \emph{diagramem elementarnym} nazywamy zbiór
\[
\operatorname{ED}(\mathcal M):=
\{\sigma:\ \sigma\text{ jest zdaniem języka }\mathcal L(M)
\text{ i }\mathcal M_M\models\sigma\}.
\]
W diagramie umieszczamy wszystkie prawdziwe zdania
rozszerzonego języka, również z negacjami
i kwantyfikatorami, a nie tylko zdania atomowe.
\end{definicja}

\begin{przyklad}[Jakie informacje zawiera diagram?]
\label{ex:ls-diagram}
Dla różnych $a,b\in M$ diagram zawiera
$c_a\ne c_b$. Jeśli $f^{\mathcal M}(a)=b$
dla jednoargumentowego symbolu $f$, zawiera
$f(c_a)=c_b$. Jeśli $\mathcal M\models
\exists y\,\theta(y,a)$, diagram zawiera zdanie
$\exists y\,\theta(y,c_a)$. Ostatni zapis
pokazuje, dlaczego diagram elementarny zawiera
więcej informacji niż tabela wartości działań i relacji.
\end{przyklad}

\begin{lemat}[Diagram daje zanurzenie elementarne]
\label{lem:ls-diagram-embedding}
Niech $\mathcal M$ będzie strukturą języka
$\mathcal L$ o niepustym uniwersum $M$.
Jeżeli struktura $\mathcal B$ języka $\mathcal L(M)$
spełnia $\operatorname{ED}(\mathcal M)$,
to funkcja $j:M\to B$ określona przez
$j(a):=c_a^{\mathcal B}$
jest zanurzeniem elementarnym
$\mathcal M$ w redukt $\mathcal B|_{\mathcal L}$.
\end{lemat}
\begin{proof}
Jeśli $a\ne b$, zdanie $c_a\ne c_b$
należy do diagramu i jest prawdziwe w $\mathcal B$,
zatem $j(a)\ne j(b)$.
Niech $f$ będzie symbolem funkcyjnym
o arności $k$, a $\vec a\in M^k$.
Diagram zawiera zdanie
$f(c_{a_1},\ldots,c_{a_k})=
c_{f^{\mathcal M}(\vec a)}$,
więc
$f^{\mathcal B}(j(\vec a))=
j(f^{\mathcal M}(\vec a))$.
Dla symboli stałych języka pierwotnego
stosujemy tę samą równość w arności $0$.
Zdania atomowe i ich negacje zapewniają
zachowanie wszystkich relacji.

Pozostaje elementarność. Dla dowolnej formuły
$\varphi(\vec x)$ języka $\mathcal L$
i krotki $\vec a\in M^{|\vec x|}$,
jeśli $\mathcal M\models\varphi(\vec a)$,
to $\varphi(c_{\vec a})\in
\operatorname{ED}(\mathcal M)$,
zatem $\mathcal B\models\varphi(j(\vec a))$.
Jeśli natomiast $\varphi(\vec a)$ jest fałszywa
w $\mathcal M$, diagram zawiera
$\neg\varphi(c_{\vec a})$,
więc $\mathcal B\not\models\varphi(j(\vec a))$.
W ten sposób dla każdej formuły i wszystkich
parametrów zachodzi równoważność wymagająca
elementarnego zanurzenia.
\end{proof}

\section{Górne twierdzenie Löwenheima--Skolema}
\begin{twierdzenie}[Górne twierdzenie Löwenheima--Skolema]
\label{thm:upward-ls}
Niech $\mathcal L$ będzie językiem pierwszego rzędu,
a $\mathcal M$ jego nieskończoną strukturą
o uniwersum $M$. Dla każdej mocy nieskończonej
$\lambda$ spełniającej
$\lambda\ge\max(|M|,|\mathcal L|,\aleph_0)$
istnieje struktura $\mathcal N$ języka $\mathcal L$
taka, że $\mathcal M\preccurlyeq\mathcal N$
oraz $|N|=\lambda$.
\end{twierdzenie}
\begin{proof}
\emph{Etap 1: wymuszenie co najmniej $\lambda$ elementów.}
W języku $\mathcal L(M)$ z poprzedniej sekcji
wprowadzamy dodatkowe, nowe stałe $d_i$,
gdzie $i$ przebiega zbiór indeksów $I$ mocy $\lambda$.
Rozważamy teorię
\[
\Gamma:=\operatorname{ED}(\mathcal M)
\cup\{d_i\ne d_j:i,j\in I,\ i\ne j\}.
\]
Każdy skończony fragment $\Gamma_0\subseteq\Gamma$
ma model. Istotnie, wszystkie stałe $c_a$
interpretujemy w $\mathcal M$ jako $a$, a skończonej
liczbie symboli $d_i$ występujących w $\Gamma_0$
przyporządkowujemy parami różne elementy
nieskończonego $M$. Pozostałym symbolom
przypisujemy ustalony element $M$.
Wtedy diagram jest spełniony, podobnie jak
wszystkie nierówności należące do $\Gamma_0$.
Twierdzenie o zwartości~\ref{thm:compactness-fol}
daje model $\mathcal B\models\Gamma$.
Ponieważ wartości $d_i^{\mathcal B}$
są parami różne, $|B|\ge\lambda$.

\emph{Etap 2: zachowanie starego modelu.}
Lemat~\ref{lem:ls-diagram-embedding}
daje zanurzenie elementarne
$j:\mathcal M\to\mathcal B|_{\mathcal L}$
określone przez $j(a)=c_a^{\mathcal B}$.
Na tym etapie $\mathcal M$ jest izomorficzna
z podstrukturą $j[M]$, ale nie musi być
dosłownie podstrukturą $\mathcal B$.

\emph{Etap 3: uzyskanie dokładnie mocy $\lambda$.}
Przechodzimy do reduktu
$\mathcal B|_{\mathcal L}$ i wybieramy
zbiór $X:=j[M]\cup\{d_i^{\mathcal B}:i\in I\}$.
Ma on moc dokładnie $\lambda$:
jest sumą zbiorów mocy co najwyżej
$|M|\le\lambda$ i $\lambda$,
a drugi z nich ma moc $\lambda$.
Z Twierdzenia~\ref{thm:downward-ls},
zastosowanego do reduktu i do $X$,
otrzymujemy $\mathcal C\preccurlyeq
\mathcal B|_{\mathcal L}$
z $X\subseteq C$ i $|C|=\lambda$.
Dla każdej formuły $\varphi(\vec x)$
języka $\mathcal L$ i krotki $\vec a\in M^{|\vec x|}$
zachodzi
\[
\mathcal M\models\varphi(\vec a)
\ \Longleftrightarrow\
\mathcal B|_{\mathcal L}\models\varphi(j(\vec a))
\ \Longleftrightarrow\
\mathcal C\models\varphi(j(\vec a)).
\]
Pierwsza równoważność wynika z diagramu,
druga z elementarności $\mathcal C$.
Także $j:\mathcal M\to\mathcal C$
jest więc zanurzeniem elementarnym.

\emph{Etap 4: przejście od kopii do dosłownego rozszerzenia.}
Wybieramy zbiór $D$ rozłączny z $M$,
równoliczny z $C\setminus j[M]$,
i bijekcję $h:M\cup D\to C$
spełniającą $h(a)=j(a)$ dla $a\in M$.
Przenosimy interpretacje symboli z $\mathcal C$
na $M\cup D$ za pomocą $h$.
Dla symbolu funkcyjnego $f$ arności $k$
przyjmujemy
$f^{\mathcal N}(\vec u)=
h^{-1}(f^{\mathcal C}(h(\vec u)))$,
a dla symbolu relacji $R$:
$R^{\mathcal N}(\vec u)\Longleftrightarrow
R^{\mathcal C}(h(\vec u))$.
Symbole stałych interpretujemy jako
$h^{-1}(c^{\mathcal C})$.
Ponieważ $j$ zachowuje wszystkie symbole,
interpretacje na $M$ pokrywają się
z pierwotnymi interpretacjami w $\mathcal M$.
Równoważność formuł z etapu 3 przenosi się
przez izomorfizm $h$, zatem
$\mathcal M\preccurlyeq\mathcal N$.
Bijekcja $h$ gwarantuje $|N|=|C|=\lambda$.
\end{proof}

\begin{wniosek}[Modele wszystkich nieskończonych mocy]
\label{cor:ls-all-infinite-cardinals}
Niech $T$ będzie zbiorem zdań pierwszego rzędu
w języku co najwyżej przeliczalnym.
Jeżeli $T$ ma choć jeden model nieskończony,
to dla każdej nieskończonej mocy $\lambda$
istnieje model $T$ o mocy $\lambda$.
\end{wniosek}
\begin{proof}
Na mocy Wniosku~\ref{cor:downward-countable}
istnieje przeliczalnie nieskończony model
$\mathcal M_0\models T$.
Dla $\lambda=\aleph_0$ jest on wystarczający.
Dla $\lambda>\aleph_0$
stosujemy Twierdzenie~\ref{thm:upward-ls}
do $\mathcal M_0$ i wybranej $\lambda$.
Rozszerzenie elementarne spełnia wszystkie
zdania prawdziwe w $\mathcal M_0$,
w szczególności wszystkie aksjomaty $T$.
\end{proof}

\begin{przyklad}[Dlaczego założenie nieskończoności jest konieczne?]
\label{ex:ls-finite-models}
Zdanie
\[
\exists x_1\cdots\exists x_5\left(
\bigwedge_{1\le i<j\le5}x_i\ne x_j
\ \land\
\forall y\bigvee_{i=1}^{5}y=x_i\right)
\]
stwierdza, że uniwersum ma dokładnie pięć elementów.
Każdy jego model jest pięcioelementowy.
Nie może mieć nieskończonego rozszerzenia elementarnego,
ponieważ zdania pierwszego rzędu są zachowywane
przez takie rozszerzenia.
\end{przyklad}

\section{Łańcuchy elementarne i ich suma}
Można też powiększać modele etapami.
Zachowanie elementarności przy przejściu do
sumy nie wynika tylko z definicji sumy:
trzeba skontrolować kwantyfikatory.

\begin{lemat}[Suma rosnącego ciągu elementarnego]
\label{lem:ls-elementary-chain}
Niech $(\mathcal M_n)_{n\in\N}$ będzie ciągiem
struktur jednego języka $\mathcal L$
takim, że $\mathcal M_n\preccurlyeq\mathcal M_{n+1}$
dla każdego $n$. Wtedy na
$M:=\bigcup_{n\in\N}M_n$ istnieje naturalna
struktura $\mathcal M$ zgodna z interpretacjami
na wszystkich etapach i dla każdego $n$
zachodzi $\mathcal M_n\preccurlyeq\mathcal M$.
\end{lemat}
\begin{proof}
Każda skończona krotka elementów $M$
należy do pewnego $M_k$: wybieramy
maksimum skończenie wielu indeksów etapów,
na których pojawiają się współrzędne.
Interpretacje stałych i funkcji na etapach
są zgodne z inkluzjami podstruktur;
relacje nie zmieniają wartości
na elementach poprzednich etapów.
Dlatego na sumie otrzymujemy strukturę
$\mathcal M$.

Elementarność wykazujemy indukcją
po budowie formuły $\varphi(\vec x)$,
jednocześnie dla każdego $n$ i
$\vec a\in M_n^{|\vec x|}$.
Dla formuł atomowych zgodność
wynika z definicji struktury na sumie,
a dla spójników z założenia indukcyjnego.
Niech $\varphi(\vec x):=\exists y\,\theta(y,\vec x)$.
Jeżeli $\mathcal M_n\models
\exists y\,\theta(y,\vec a)$,
istnieje $b\in M_n$ z
$\mathcal M_n\models\theta(b,\vec a)$,
więc indukcja daje
$\mathcal M\models\theta(b,\vec a)$.
Odwrotnie, gdy
$\mathcal M\models\exists y\,\theta(y,\vec a)$,
wybieramy świadka $b\in M$, a następnie
$k\ge n$ z $b\in M_k$.
Założenie indukcyjne daje
$\mathcal M_k\models\theta(b,\vec a)$,
więc $\mathcal M_k\models
\exists y\,\theta(y,\vec a)$.
Przechodniość elementarności na
skończonym łańcuchu od $n$ do $k$
daje $\mathcal M_n\models
\exists y\,\theta(y,\vec a)$.
Kwantyfikator ogólny sprowadzamy
do negacji i egzystencjalnego.
To kończy indukcję.
\end{proof}

\begin{uwaga}[Dłuższe łańcuchy]
\label{rem:ls-transfinite-chains}
Analogiczne twierdzenie obowiązuje
dla odpowiednio skierowanych rodzin
i dla łańcuchów indeksowanych
liczbami porządkowymi. Dokładniejsze
sformułowanie wymaga pojęć, które
wprowadzimy w części teoriomnogościowej.
Dla górnego twierdzenia
Löwenheima--Skolema wystarczyła nam
bezpośrednia konstrukcja z diagramem i zwartością.
\end{uwaga}

\section{Niestandardowe modele prawdziwej arytmetyki}
\begin{twierdzenie}[Przeliczalny niestandardowy model prawdziwej arytmetyki]
\label{thm:ls-true-arithmetic-nonstandard}
Niech $\operatorname{Th}(\mathfrak N)$
będzie zbiorem wszystkich zdań języka
arytmetyki pierwszego rzędu prawdziwych
w standardowej strukturze $\mathfrak N$.
Istnieje przeliczalnie nieskończona
struktura arytmetyczna $\mathcal M$
spełniająca $\operatorname{Th}(\mathfrak N)$
oraz zawierająca element większy
od interpretacji każdego standardowego
numerału. Struktura ta nie jest
izomorficzna ze $\mathfrak N$.
\end{twierdzenie}
\begin{proof}
Do języka arytmetyki dodajemy nową
stałą $c$ i teorię
\[
\Gamma:=
\operatorname{Th}(\mathfrak N)
\cup\{\overline n<c:n\in\N\}.
\]
Każdy skończony podzbiór $\Gamma_0$
ma model: interpretujemy arytmetykę
standardowo, a $c$ jako liczbę naturalną
większą od wszystkich wartości $n$,
dla których zdanie $\overline n<c$
występuje w $\Gamma_0$.
Z Twierdzenia o zwartości~\ref{thm:compactness-fol}
wynika model $\mathcal B\models\Gamma$.
Rozszerzony język jest przeliczalny,
a model $\mathcal B$ jest nieskończony,
ponieważ spełnia prawdziwe zdania
arytmetyczne stwierdzające istnienie
dowolnie wielu różnych elementów.
Wniosek~\ref{cor:downward-countable},
zastosowany do języka z dodatkową
stałą $c$, daje przeliczalnie nieskończony
$\mathcal M_c\preccurlyeq\mathcal B$.
Zawiera on interpretację $c$,
ponieważ jest podstrukturą w tym języku.
Redukt $\mathcal M$ do języka arytmetyki
spełnia $\operatorname{Th}(\mathfrak N)$
oraz ma przeliczalny uniwersum.

Gdyby $\mathcal M$ była izomorficzna
ze $\mathfrak N$, każdy element $M$
byłby obrazem pewnego standardowego
numerału: izomorfizm zachowuje $0$
i następnik $S$. W szczególności
$c^{\mathcal M_c}=\overline k^{\mathcal M}$
dla pewnego $k\in\N$. Tymczasem
$\mathcal M_c\models\overline k<c$.
Po podstawieniu otrzymalibyśmy
$\mathcal M\models\overline k<\overline k$,
wbrew prawdziwym aksjomatom
porządku zawartym w
$\operatorname{Th}(\mathfrak N)$.
Zatem $\mathcal M$ jest niestandardowa.
\end{proof}

\begin{uwaga}[Dlaczego indukcja Peana tego nie wyklucza?]
\label{rem:ls-nonstandard-cut}
Zewnętrzny zbiór interpretacji
standardowych numerałów
$C:=\{\overline n^{\mathcal M}:n\in\N\}$
zawiera $0^{\mathcal M}$ i jest
zamknięty względem następnika.
Nie jest jednak definiowalny
w $\mathcal M$ żadną formułą
arytmetyki, nawet z parametrami
z $M$: gdyby $C$ był tak
definiowalny, odpowiednia instancja
indukcji Peana wymusiłaby $C=M$.
Jest to niemożliwe, ponieważ
$c^{\mathcal M_c}\notin C$.
Indukcja pierwszego rzędu nie
kwantyfikuje po wszystkich
zewnętrznych podzbiorach modelu.
\end{uwaga}

\section{Paradoks Skolema i różnica perspektyw}
\begin{uwaga}[Zapowiedź aksjomatów ZFC]
\label{rem:ls-zfc-intro}
ZFC jest teorią pierwszego rzędu
w języku z relacją $\in$.
Jej szczegółowe aksjomaty zostaną
wprowadzone w następnym rozdziale.
W tej sekcji wykorzystamy dwa
własne fakty ZFC: w każdym modelu
istnieją obiekty odpowiadające
dowolnie wielu różnym liczbom
naturalnym, a wewnętrzny zbiór
liczb rzeczywistych jest w ZFC
dowodliwie nieprzeliczalny.
Ten drugi fakt wynika z twierdzenia
Cantora o zbiorze potęgowym.
\end{uwaga}

\begin{twierdzenie}[Warunkowe istnienie przeliczalnego modelu ZFC]
\label{thm:skolem-paradox}
Jeżeli ZFC ma model, to ma model
o przeliczalnie nieskończonym uniwersum.
W takim modelu może być prawdziwe
zdanie stwierdzające nieprzeliczalność
jego zbioru liczb rzeczywistych.
Nie prowadzi to do sprzeczności
z nieprzeliczalnością liczb rzeczywistych.
\end{twierdzenie}
\begin{proof}
Język ZFC zawiera tylko symbol $\in$
i równość logiczną, a więc jest
co najwyżej przeliczalny.
Każdy model ZFC jest nieskończony:
aksjomaty ZFC dowodzą istnienia
obiektów odpowiadających dowolnie
wielu różnym liczbom naturalnym.
Jeśli $\mathcal B\models\mathrm{ZFC}$,
to Wniosek~\ref{cor:downward-countable}
daje przeliczalnie nieskończoną
podstrukturę elementarną
$\mathcal S\preccurlyeq\mathcal B$,
która również spełnia wszystkie
aksjomaty ZFC. W szczególności
$\mathcal S$ spełnia zdanie
o nieprzeliczalności własnego
wewnętrznego zbioru liczb rzeczywistych.

Niech $S$ będzie uniwersum
$\mathcal S$, a $r\in S$ elementem,
który $\mathcal S$ interpretuje jako
zbiór liczb rzeczywistych.
Zewnętrznie rozpatrywany zbiór
\[
E_r:=\{a\in S:\mathcal S\models a\in r\}
\]
jest co najwyżej przeliczalny,
gdyż $E_r\subseteq S$.
Nie znaczy to, że istnieje
\emph{element $S$}, który model
rozpoznaje jako bijekcję między
jego wewnętrznym zbiorem liczb
naturalnych a $r$.
Nieprzeliczalność wewnętrzna
stwierdza właśnie brak takiego
elementu wewnątrz modelu.
Zewnętrzne wyliczenie $E_r$
nie musi być reprezentowane
przez żaden obiekt $\mathcal S$,
a ponadto wewnętrzny zbiór liczb
naturalnych modelu nie musi być
standardowy. Obie wypowiedzi
mają różne zakresy kwantyfikatorów.
\end{proof}

\begin{uwaga}[Zbiór potęgowy jest również wewnętrzny]
\label{rem:skolem-powerset}
Aksjomat zbioru potęgowego w ZFC
zapewnia istnienie obiektu
zawierającego wszystkie podzbiory
danego zbioru w znaczeniu modelu.
Kwantyfikatory przebiegają jednak
po elementach modelu, a nie po
wszystkich zewnętrznych podzbiorach
jego uniwersum. Dlatego model
nie musi zawierać wszystkich
zewnętrznych podzbiorów
$E_x:=\{a\in S:\mathcal S\models a\in x\}$.
\end{uwaga}

\section{Logika drugiego rzędu i granice uzyskanych wyników}
Definicje składni oraz semantyki pełnej
i semantyki Henkina wprowadziliśmy
w rozdziale~\ref{ch:semantics-fol},
w sekcji~\ref{sec:second-order-introduction}.
W pełnej semantyce kwantyfikator
po zmiennej jednoargumentowej relacji
przebiega po wszystkich \emph{zewnętrznych}
podzbiorach uniwersum.

\begin{twierdzenie}[Kategoryczność arytmetyki Peana drugiego rzędu]
\label{thm:second-order-peano-categorical}
Niech $\mathcal M$ będzie strukturą
spełniającą standardowe aksjomaty
Peana drugiego rzędu w pełnej
semantyce, w której zmienne
jednoargumentowych relacji
przebiegają po wszystkich
zewnętrznych podzbiorach uniwersum.
Wtedy $\mathcal M$ jest izomorficzna
ze standardową strukturą arytmetyczną
$\mathfrak N$.
\end{twierdzenie}
\begin{proof}
W metateorii definiujemy
$e:\N\to M$ przez
$e(n):=(S^{\mathcal M})^n(0^{\mathcal M})$.
Aksjomaty Peana zapewniają, że
$S^{\mathcal M}$ jest iniekcją
i $0^{\mathcal M}$ nie należy
do obrazu następnika.
Jeśli $e(m)=e(n)$ dla $m<n$,
skracamy $m$-krotne zastosowanie
iniekcji $S^{\mathcal M}$,
otrzymując
$0^{\mathcal M}=
(S^{\mathcal M})^{n-m}(0^{\mathcal M})$,
sprzecznie z aksjomatem o zerze.
Zatem $e$ jest iniekcją.

Niech $E:=\{e(n):n\in\N\}$.
Zawiera $0^{\mathcal M}$
i jest zamknięty względem
$S^{\mathcal M}$. Pełna semantyka
drugiego rzędu pozwala zastosować
aksjomat indukcji \emph{do tego właśnie}
zewnętrznie określonego zbioru $E$.
Otrzymujemy $E=M$,
czyli surjektywność $e$.
Równania rekurencyjne Peana dla
dodawania i mnożenia zapewniają
przez indukcję po drugim argumencie,
że $e(m+n)=e(m)+^{\mathcal M}e(n)$
i $e(mn)=e(m)\cdot^{\mathcal M}e(n)$.
Relacja $<$, jeśli należy do
sygnatury, jest wyznaczona
przez działania i relację
następnika zgodnie z aksjomatami.
Tak więc $e$ jest izomorfizmem
struktur arytmetycznych.
\end{proof}

\begin{uwaga}[Pełna semantyka a semantyka Henkina]
\label{rem:ls-second-order-scope}
W semantyce Henkina zmienne drugiego
rzędu przebiegają tylko po zadanej
rodzinie podzbiorów uniwersum.
Zewnętrznie określony zbiór
$E=\{e(n):n\in\N\}$ nie musi
należeć do tej rodziny;
dowód surjektywności $e$
nie przenosi się więc automatycznie
na semantykę Henkina. Kategoryczność
drugiego rzędu przy pełnej semantyce
nie przeczy twierdzeniom
Löwenheima--Skolema dla pierwszego rzędu.
\end{uwaga}

\begin{uwaga}[Niezupełność a mnogość modeli]
\label{rem:ls-incompleteness-distinction}
Brak dowodu zdania $\varphi$ w teorii
nie oznacza, że $\varphi$ jest fałszywe.
Niezupełność efektywnej teorii arytmetycznej
i istnienie nieizomorficznych modeli
jednej teorii pierwszego rzędu to różne wyniki.
Teoria zupełna pierwszego rzędu może mieć
modele różnych mocy, a więc nieizomorficzne,
mimo że rozstrzyga każde zdanie swego języka.
Oddzielne pozostają także definiowalność,
efektywna wyliczalność i rozstrzygalność.
\end{uwaga}
